\begin{table*}[t]
\centering
\caption{RULER~\citep{hsieh2024ruler} 与 LongBench-v2~\citep{bai2024longbench2} 上的模型性能。我们报告平均性能（Perf.）与 $\Omega_{\mathrm{{MSR}}}$。}
\small
\resizebox{\linewidth}{!}{
\begin{tabular}{l | c c c c c c | c | c | c c | c | c}
\toprule
\multirow{2.5}{*}{\textbf{模型}} & \multicolumn{8}{c|}{\textbf{\texttt{RULER}}} & \multicolumn{4}{c}{\textbf{\texttt{LongBench-v2}}} \\
\cmidrule(lr){2-9} \cmidrule(lr){10-13}
 & \bf 8K &\bf 16K &\bf 32K &\bf 64K &\bf 128K &\bf 256K & \textbf{性能} & $\mathbf{\Omega_{\mathrm{\textbf{MSR}}}}$ &\bf Easy &\bf Hard & \textbf{性能} & $\mathbf{\Omega_{\mathrm{\textbf{MSR}}}}$ \\
\arrayrulecolor{black}\midrule

% ==========================================
% SECTION 1: Qwen3-4B (Light Blue)
% ==========================================
\rowcolor{blue!10}
\multicolumn{13}{c}{\textbf{Qwen3-4B 骨干模型}} \\
\arrayrulecolor{black!20}\midrule
Qwen3-4B & 87.49 & 86.82 & 60.05 & 70.98 & 53.19 & 43.27 & 66.00 & - & 32.67 & 22.18 & 25.96 & - \\
+ InfLLM-V2 & 78.59 & 74.40 & 43.39 & 44.94 & 28.57 & 27.01 & 51.02 & - & 28.00 & 24.06 & 25.48 & - \\
+ DuoAttention & 76.82 & 75.17 & 50.53 & 61.42 & 45.86 & 45.46 & 58.30 & 0.70 & 28.67 & \underline{24.44} & 25.96 & 0.70\\
+ PruLong & 77.03 & 73.99 & \underline{51.32} & 60.98 & 42.70 & \textbf{48.76} & 58.38& 0.70 & 27.60 & 22.33 & 24.35 & 0.70\\
\arrayrulecolor{black!40}\midrule
+ Elastic Attention (FA-SSA) & \underline{83.35} & \underline{79.24} & 50.79 & \underline{67.03} & \underline{47.83} & \underline{47.32} & \underline{61.81}& 0.66 & \textbf{34.00} & \underline{24.44} & \underline{27.88} & 0.70\\
+ Elastic Attention (FA-XA) & \textbf{86.56} & \textbf{85.38} & \textbf{56.88} & \textbf{69.42} & \textbf{49.48} & 43.47 & \textbf{63.27}& 0.67 & \underline{32.00} & \textbf{25.94} & \textbf{28.12} & 0.72\\

% ==========================================
% SECTION 2: Qwen3-8B (Deeper Blue)
% ==========================================
\arrayrulecolor{black}\midrule
\rowcolor{blue!25}
\multicolumn{13}{c}{\textbf{Qwen3-8B 骨干模型}} \\
\arrayrulecolor{black!20}\midrule
Qwen3-8B & 89.69 & 85.62 & 63.23 & 82.39 & 65.84 & 66.71 & 75.74 & - & 39.33 & 27.82 & 31.97 & - \\
+ InfLLM-V2 & 77.01 & 68.53 & 25.93 & 53.47 & 34.96 & 32.95 & 52.58 & - & 29.32 & 28.67 & 29.09 & -\\
+ DuoAttention & 81.03 & 78.85 & 56.61 & 70.10 & 54.68 & 56.28 & 65.94 & 0.70 & \textbf{40.67} & 25.56 & 31.01 & 0.70\\
+ PruLong & \textbf{87.20} & 82.69 & 61.64 & 73.05 & 59.74 & 58.82 & 69.90 & 0.70 & \underline{38.67} & 27.82 & \underline{31.73} & 0.70\\
\arrayrulecolor{black!40}\midrule
+ Elastic Attention (FA-SSA) & \underline{86.62} & \underline{82.81} & \underline{64.55} & \underline{77.41} & \underline{61.17} & \underline{61.75} & \underline{71.74} & 0.65 & 37.33 & \textbf{31.20} & \textbf{33.41} & 0.66\\
+ Elastic Attention (FA-XA) & 85.07 & \textbf{85.12} & \textbf{65.08} & \textbf{82.34} & \textbf{64.57} & \textbf{63.41} & \textbf{73.87} & 0.76 & 30.68 & \underline{30.77} & 30.74 & 0.78\\

% ==========================================
% SECTION 3: Llama-3.1-8B (Orange)
% ==========================================
\arrayrulecolor{black}\midrule
\rowcolor{cyan!20}
\multicolumn{13}{c}{\textbf{Llama-3.1-8B-Instruct 骨干模型}} \\
\arrayrulecolor{black!20}\midrule
Llama-3.1-8B-Instruct & 92.88 & 92.83 & 89.46 & 70.79 & 80.12 & 72.34 & 83.47 & - & 32.00 & 33.08 & 32.69 & - \\
+ InfLLM-V2 & 89.30 & 80.93 & 60.98 & 35.90 & 32.29 & 47.27 & 59.10 & - & \underline{29.32} & 28.67 & \underline{29.09} & - \\
+ DuoAttention & 87.20 & 77.85 & 66.99 & 40.13 & 57.45 & 42.92 & 62.92 & 0.70 & 28.67 & 27.44 & 27.88 & 0.70 \\
+ PruLong & 83.54 & 69.35 & 56.83 & 30.88 & 33.74 & 21.64 & 48.82 & 0.70 & 28.00 & 28.92 & 28.61 & 0.70\\
\arrayrulecolor{black!40}\midrule
+ Elastic Attention (FA-SSA) & \underline{89.93} & \underline{83.42} & \underline{80.20} & \underline{56.30} & \underline{68.16} & \underline{56.47} & \underline{72.85} & 0.65 & 28.00 & \underline{29.70} & \underline{29.09} & 0.68\\
+ Elastic Attention (FA-XA) & \textbf{92.82} & \textbf{92.00} & \textbf{87.80} & \textbf{68.23} & \textbf{78.87} & \textbf{68.51} & \textbf{81.82} & 0.72 & \textbf{30.67} & \textbf{30.83} & \textbf{30.77} & 0.75 \\
\arrayrulecolor{black}\bottomrule
\end{tabular}
}
\vspace{-0.5em}
\label{tab:ruler_lb2_result}
\end{table*}


\section{实验}
\label{sec:experiment}

\subsection{实验设置}
\label{subsec:exp_settings}

\paragraph{训练与数据}
\label{subsec:train_settings}
我们选择 Qwen3-(4B/8B)~\citep{yang2025qwen3technicalreport} 和 Llama-3.1-8B-Instruct~\citep{grattafiori2024llama} 作为骨干 LLM。
训练数据由 5 个来源构建：ChatQA2-Long-SFT-data~\citep{xu2024chatqa}、MuSiQue~\citep{trivedi2022musique}、CoLT-132K~\citep{li2025aixcoder}、GovReport~\citep{huang-etal-2021-efficient} 与 XSum~\citep{xsum-emnlp}，覆盖稀疏敏感任务（单文档 QA、多跳 QA）与稀疏鲁棒任务（代码补全、摘要、上下文学习）。
最终数据集序列长度范围为 8K 到 64K token，总计约 0.74B token。
针对稀疏鲁棒与稀疏敏感任务，我们经验设置目标稀疏度分别为 $\boldsymbol{t}=1.0$ 与 $\boldsymbol{t}=0.7$。
我们在 $8\times$A800 GPU 上训练，每次实验可在 12 小时内完成。
更多训练细节见附录~\ref{appendix:implementation_details}，超参数见表~\ref{tab:combined_hyperparameters}。

\vspace{-0.5em}
\paragraph{评测}
\label{subsec:baseline}
我们将方法与代表性训练式稀疏方法比较，包括 DuoAttention~\citep{xiaoduoattention2025}、PruLong~\citep{bhaskar2025cache} 与 InfLLM-V2~\citep{infllmv2}。
对于稀疏头的注意力计算模式，我们考虑 Streaming Sparse Attention（SSA）~\citep{streamingllm} 与 XAttention（XA）~\citep{xattention}。
在 XA 中设阈值 $\tau=0.9$，其余超参数采用默认值。
我们使用“\{检索头模式\}–\{稀疏头模式\}”表示不同头配置，例如 FA-SSA 表示检索头采用 FA、稀疏头采用 SSA。
由于篇幅限制，MoBA~\citep{moba} 与 NSA~\citep{NSA} 的实验，以及更详细基线配置见附录~\ref{appdix:eval_res}。
全部评测均在 \textbf{\texttt{LOOM-Eval}} 框架~\citep{tang2025loom} 上完成。

\subsection{评测结果}
\label{subsec:eval_res_main}
\paragraph{真实长上下文任务}
我们首先在 LongBench-E~\citep{bai2024longbench} 上评估方法。该基准包含 6 类共 14 个真实长上下文任务。
表~\ref{tab:longbench_main} 对比了 Elastic Attention 与其他基线。
在每个比较组中，所有方法使用相同骨干模型与相同训练数据。
结果显示，我们的方法在平均性能上持续最优，能够在提升推理效率的同时达到或逼近骨干模型（全 FA）表现。
在不同任务上，Elastic Attention 能动态分配不同稀疏水平，例如在代码任务上平均 $\Omega_{\mathrm{MSR}}$ 约为 0.85，而在 QA 上约为 0.68。
同时我们注意到，在稀疏鲁棒任务（如 Code、Summ）上，本方法有时弱于部分基线。
其原因是我们采用更高稀疏度，而 InfLLM-V2 在部分场景会使用 FA\footnote{InfLLM-V2 没有严格定义的 $\Omega_{\mathrm{MSR}}$，因为其在上下文长度小于 8K 时使用 FA，超过 8K 时切换为混合注意力模式。}。

\paragraph{长度外推能力测试}
我们在 RULER~\citep{hsieh2024ruler} 上评测长度外推能力，考察不同上下文长度区间下模型表现。
训练时最大上下文长度为 64K token，评测长度覆盖 8K 至 256K token。
如表~\ref{tab:ruler_lb2_result}（左半部分）所示，我们的方法在所有长度上均取得最佳性能，同时保持较优稀疏水平，$\Omega_{\mathrm{MSR}}$ 基本稳定在 $0.65\sim 0.7$。
此外我们观察到，对于 8B 规模模型，当上下文超过 64K token 时，FA–XA 配置通常表现最好。
原因在于 FA–XA 的 $\Omega_{\mathrm{ESR}}$ 低于其他配置，可保留并利用更多有效信息。
这种优势在超长上下文下尤为明显，例如在 256K token 时 FA–XA 仍保持较强性能。

\paragraph{长程推理任务}
我们进一步在 LongBench-V2~\citep{bai2024longbench2} 上评测，该基准被广泛用于长上下文推理，长度范围约为 8K 至 2M 词。
如表~\ref{tab:ruler_lb2_result}（右半部分）所示，我们的方法在 Easy 与 Hard 设置下都取得了稳定强结果，并获得最佳平均性能（FA-SSA 或 FA-XA 配置）。
值得注意的是，FA–XA 在 Qwen3-8B 上未表现出同等优势，这可能与该模型自身特性有关，后续可通过更细粒度超参数调优进一步改进。



% We employ a comprehensive evaluation framework~\citep{tang2025loom} to assess the model's capabilities across three complementary dimensions. We primarily utilize \textbf{LongBench}~\citep{bai2024longbench} to evaluate real-world understanding; its diverse task set serves as a critical diagnostic tool to verify the router's ability to dynamically assign task-specific sparsity. To assess the effective context window, we employ \textbf{RULER}~\citep{hsieh2024ruler}. Furthermore, we test the model with \textbf{LongBench-v2}~\citep{bai2024longbench2}, focusing on multi-hop reasoning to demonstrate robustness in capturing intricate dependencies that standard benchmarks may overlook.
